\section{Evacuation, columns and obstacle detours}\label{sec:solid-routing}
The following refinements choose the occupied regions and the storage order
explicitly. Ordered low extraction controls a height observer; simultaneous
lifting controls a first-coordinate observer; expanded centers permit
three-dimensional detours. Fixed columns and recorded pointers give further
ways to keep a whole transition outside a detector. Each proof specifies
which parts of the geometry act on solid neighborhoods and which bounds
are needed only for the coded central plane.

Unless an individual theorem specifies another class, whole-space
comparison means smooth $u\in CH^2\cap C^1L^2$, bounded $u,\nabla u$
on finite intervals, and a pressure representative in $CH^1$, as in
Section~\ref{bc:analytic-section}. The torus comparison class is the
classical class of Lemma~\ref{lem:b-comparison}. The distinct pressure
conditions stated below remain part of their respective results.

\subsection{Ordered extraction below the computation}\label{bc:ordered-section}
\begin{proposition}[Ordered full-box transport]\label{bc:ordered}
Let $D_i,S_i$ be finite families of closed rational axis-aligned boxes in $\R^3$
with positive side lengths, pairwise disjoint within each family.
Suppose the assigned center-to-center affine map $F_i$ satisfies
$F_i(D_i)=S_i$ and has positive diagonal linear part $A_i$ of
determinant one. There is an effective compactly supported solenoidal
unit-time field, zero near the time endpoints, whose time-one map has
this action on every full $D_i$. If both centers of a pair have negative
third coordinate, every point of $D_i$ whose third coordinate equals that
of its center follows a path with negative third coordinate.
\end{proposition}
\begin{proof}\label{bc:ordered-proof}
Choose an integer $L>1$ larger than the absolute values of all endpoint
coordinates. The storage columns have first coordinates $U_i=4(i+1)L$
and common height $Z=-4L$. Write $[\ell_i,r_i]$ for the first-coordinate
interval of $D_i$. Extract sources in decreasing order of $r_i$.
Translate each horizontally right to its column and then vertically to
$Z$, retaining its second coordinate. If an unextracted box meets the
active box in both the second and third projections, disjointness forces
their first intervals to be separated. That box cannot lie to the right:
its right endpoint would then exceed $r_i$, contrary to the extraction
order. It lies strictly to the left of $\ell_i$ and misses the entire
rightward sweep. Stored boxes lie below the original boxes; vertical
moves occur in distinct columns outside them.

At storage height, adjust the second center coordinate to its target
value and apply $A_i^{\theta}$. Each intermediate half-width lies between
its two endpoint values, hence is less than $L$. The columns remain
disjoint throughout these operations. Deliver boxes in increasing order
of the right endpoint of their target first interval: raise a box to its
target height and translate it left into its target. The raising column
is empty. As above, a previously delivered box whose last two intervals
meet the active ones must be separated in the first coordinate. Its
smaller right endpoint puts it wholly to the left of the active target
interval, so the leftward sweep stops before reaching it.
Undelivered boxes remain at storage height. These observations give
positive rational clearances for all swept bounding boxes. Choose cutoff
padding below each clearance and use Lemma~\ref{lem:bb-curl} in successive
flat slots. Induction over the slots proves the full-box assertion.
A point initially in the central third-coordinate plane retains zero
third-coordinate offset during each diagonal deformation. Its height is
interpolated only between the two endpoint heights and $Z$. If those
endpoint heights are negative, every intermediate height is negative.
\end{proof}

The onto hypothesis in the preceding proposition is a useful exact-box
specialization. Its target-container version, including the effective
input conditions, is given in Section~\ref{bc:contained-targets}.

\begin{theorem}[A positive-height observation]\label{bc:positive-height}
On $\R^3$, the fixed label $0$ and observation $\{x_3>0\}$ realize
halting effectively by a compactly supported general force with zero
initial velocity and pressure, bounded mixed derivatives, and period one
after time one. The force has a separate alternative in
$L^2(0,\infty;H^j)$ for every $j$, tending to zero in every spatial
supremum norm at rate $O((1+t)^{-1})$.
\end{theorem}
\begin{proof}\label{bc:positive-height-proof}
Use Lemma~\ref{lem:bb-movefirst} and odd digits in base $2m+2$. Give every
simulator state a distinct positive integer index $i_s$ and put its
third-coordinate center at $z_s=i_s$ for ready halting states and
$z_s=-i_s$ otherwise. Thicken the full rectangles of
Lemma~\ref{lem:bb-onecell} by $[z_s-1/4,z_s+1/4]$. The assigned third
scale is one, so the maps satisfy Proposition~\ref{bc:ordered}.
The rational initial code $Y_0$ is loaded from the origin along the
straight segment, by a localized translation during $[0,1]$. Repeat the
unit-time mechanism after time one. At time $1+k$ the particle is exactly
the code after $k$ defined simulator transitions. If the initial machine
halts, the loading endpoint already has positive height. If it never
halts, loading has nonpositive height and all later transitions have
negative endpoint centers; the proposition excludes a positive height
at every intervening time. Conversely, a halting checkpoint has positive
height. Lemma~\ref{bc:analytic} gives the force and uniqueness, and
Lemma~\ref{ba:log-clock} gives the separate temporal alternative.
\end{proof}

\subsection{Simultaneous lifting with a convex-coordinate guard}
\label{bc:simultaneous-section}
\begin{proposition}[Simultaneous full-box transport]\label{bc:simultaneous}
Suppose finite planar source and target rectangle families $S_i,T_i$
have positive side lengths and endpoints on the grid $B^{-2}\Z^2$,
with integer $B\ge2$, and are separately
disjoint. Require their assigned center-to-center maps to satisfy
$F_i(S_i)=T_i$ and to have linear part $\diag(\lambda_i,\lambda_i^{-1})$ with
$\lambda_i\in\{B^{-1},1,B\}$. There is an effective solenoidal motion
of the full boxes obtained by taking the product with $[-1,1]$.
For every transported point its first coordinate is a convex combination
of its initial and final first coordinates at every time.
\end{proposition}
\begin{proof}\label{bc:simultaneous-proof}
Let $p_i^-,p_i^+$ be the centers in the plane $z=0$, set $h_i=5i$, and
choose a smooth $\Theta$ equal to zero on $(-\infty,1/4]$ and one on
$[3/4,\infty)$. For $0\le t\le1$, put
\[
 \theta=\Theta(3t-1),\quad l_i=1-\theta+\theta\lambda_i,
 \quad c_i=(1-\theta)p_i^-+\theta p_i^+
       +h_i\{\Theta(3t)-\Theta(3t-2)\}e_3.
\]
The desired path is $g_i(t,y)=c_i(t)+\diag(l_i,l_i^{-1},1)(y-p_i^-)$.
Use a source padding $\varepsilon=1/(10B^3)$, transported by this same
affine map, for the cutoff neighborhood. During lifting the planar
source gap is at least $B^{-2}$. During lowering the enlarged target
padding is at most $B\varepsilon<B^{-2}/2$. During the middle phase
the height difference is at least five, whereas the sum of third
half-widths is $2+2\varepsilon<5$. Thus these moving cutoff supports
can be chosen disjoint. Sum the localized curls of
Lemma~\ref{lem:bb-curl}. Differentiation and uniqueness give the displayed
paths for the full boxes. Finally
\[
 g_{i,1}(t,y)=(1-\theta)y_1+
 \theta\{p_{i,1}^++\lambda_i(y_1-p_{i,1}^-)\},
\]
which proves the asserted guard. Flat time collars hold automatically.
\end{proof}

\begin{theorem}[A zero-blank half-space test]\label{bc:zero-blank}
The fixed label $0$ and event $\{x_1<-1\}$ in $\R^3$ realize halting
with the force and comparison conclusions of Theorem~\ref{bc:positive-height}.
This construction uses a zero joint-blank digit, so every initial tape
coordinate is a finite radix sum. Its temporal alternative belongs to
$L^2(0,\infty;H^j)$ for every $j$.
\end{theorem}
\begin{proof}\label{bc:zero-blank-proof}
Use the frontier recorder, zero-blank base $B=2m$, and state first-coordinate
offsets $3j$ for nonterminal controls and $-3j$ for terminal controls,
with positive distinct indices. Encode a state by $(o_s+L,R,0)$.
The full branches are on the $B^{-2}$ grid and satisfy the preceding
proposition. Load the rational initial code from zero along a straight
segment. Nonhalting codes have nonnegative first coordinate; the whole
segment and every nonhalting transition therefore avoid $x_1<-1$.
Terminal codes have first coordinate at most $-2$, including an initially
terminal loading endpoint. Exact integer-time simulation gives both
implications. The force conclusions follow from the common analytic and
logarithmic-clock lemmas, with the compact support needed for all $H^j$
norms. No claim of periodicity is made for the slowed alternative.
\end{proof}

\subsection{Expanded centers and obstacle detours}\label{bc:expanded-height-section}
\begin{proposition}[Full solid boxes with a height guard]
\label{bc:expanded-height-routing}
Let the source and target families be finite collections of full rational
axis-aligned solid boxes, each family pairwise disjoint. Write their centers
as $a_i,b_i$ and their positive half-widths as $r_{i,j},w_{i,j}$.
Require each prescribed affine map to send its entire source box onto
its specified target box:
\[
 F_i(a_i+h)=b_i+D_i h,\qquad
 D_i=\diag(w_{i,1}/r_{i,1},w_{i,2}/r_{i,2},w_{i,3}/r_{i,3}),
 \qquad \det D_i=1.
\]
Suppose every source center has height at most $-1$, and every target
center has height at most $-1$ or exactly $1$. They admit an effective
compactly supported solenoidal unit-time realization. If a target center
has negative height, every point $a_i+h$ of its source box with $h_3=0$
remains at negative height throughout. A rational bound for the
entire support is computable from the boxes.
\end{proposition}
\begin{proof}\label{bc:expanded-height-routing-proof}
The empty family is realized by zero. Otherwise let $R=2+$ the largest
endpoint half-width, and choose an integer $H>20R$ so large that the
expanded source centers $A_i^*=Ha_i$ and target centers $B_i^*=Hb_i$
are separately more than $20R$ apart in the sup norm. Set
\[
 Z_0=1+100R+\max_i\{\|A_i^*\|_\infty,\|B_i^*\|_\infty\},
 \qquad P_i=(0,0,-Z_0-100Ri).
\]
Each of the unions $\{A_i^*\}\cup\{P_i\}$ and
$\{B_i^*\}\cup\{P_i\}$ has separation greater than $20R$.
The construction has five phases: expand source centers by a factor
from one to $H$ while keeping shapes fixed; move them successively to
$P_i$; change each parked shape by the prescribed positive diagonal
power; move them successively to $B_i^*$; contract target centers to
$b_i$ while keeping the final shapes fixed.

The simultaneous expansion is collision free. For each pair of source
boxes some coordinate has positive gap
$|a_{i,j}-a_{k,j}|-r_{i,j}-r_{k,j}$; multiplication of the center
difference by a factor at least one preserves that gap. The same
argument applies to contraction toward the disjoint targets. Choose
padding less than one quarter of the minimum such gap, and less than
one. During a parked deformation the $j$th half-width is
$r_{i,j}^{1-\theta}w_{i,j}^{\theta}\le\max(r_{i,j},w_{i,j})<R$.
Thus parked deformations have all half-widths less than $R$ and
supports of half-width at most $2R$; the storage separation suffices.

For a successive transfer, surround each stationary center by its closed
sup-norm cube of radius $4R$. These cubes are pairwise disjoint. Start
with the straight segment between the moving centers. Each nontrivial
intersection with such a cube can be replaced by a polygonal path on
its boundary: join the entrance to a vertex of its face, follow cube
edges to an exit-face vertex, and join to the exit. Fixed index orders
choose the vertices and edges. All coordinates and intersection parameters
are rational. Tangencies need no detour. Since different obstacle cubes
are disjoint, the resulting finite polygon stays at distance at least
$4R$ from every stationary center. Move along its segments in flat slots,
using a cutoff support of half-width at most $2R$. Stationary boxes have
half-width less than $R$, so no such support touches them. Lemma~\ref{lem:bb-curl}
realizes each translation and deformation.

For the height assertion, a transfer with negative endpoints has endpoints
of height at most $-H$. Its straight segment is negative. It cannot meet
an obstacle centered at positive height $H$, whose bottom is $H-4R>0$.
Every boundary used around a negative-height obstacle lies below
$-H+4R<0$. Storage heights are negative as well. Expansion, contraction,
and positive diagonal deformation preserve the central-plane condition.
Thus all these paths are negative. Finitely many rational vertices,
endpoint widths, and cutoff paddings give a rational support bound.
\end{proof}

\begin{theorem}[A height test inside one torus chart]\label{bc:expanded-height-event}
On the unit torus, the fixed label $o=(1/2,1/2,1/2)$ and event
\begin{equation}\label{bc:expanded-height-fixed-event}
 \{x:x_3\pmod1\in(1/2,1)\}
\end{equation}
realize halting by a smooth general mean-zero force, periodic after time
one, with zero initial velocity and pressure and bounded mixed derivatives.
The repeated mechanism depends on the machine alone; only initialization
depends on its finite input. The solution is unique in the torus class.
\end{theorem}
\begin{proof}\label{bc:expanded-height-event-proof}
Use the frontier recorder, odd base $B=2m+1$, distinct negative integer
heights for nonterminal states and height one for the terminal ready state.
Thicken every branch by a third half-width $1/4$. The reciprocal tape
map is onto its target rectangle; adjoining the unchanged third interval
gives an onto box map with linear part $\diag(B^{-e},B^e,1)$ and
determinant one. All source centers are
negative because terminal controls have no outgoing rules. Apply the
preceding proposition. The logical initial code $Y_0$ is rational and has
a uniform machine-dependent bound as the input varies: its two stack
coordinates lie in $[0,1]$ and its height is fixed by the initial control.
A moving translation loads the logical origin to $Y_0$, with a fixed
source half-width one and padding one. This gives a machine-dependent
rational bound $K$ for the supports of all such loaders and the repeated
mechanism; for example take $K=1+2R+M$ with $M$ an upper bound for all
routing vertices and loading coordinates and $R$ enlarged to cover their
half-widths. Choose $\lambda=1/(4K)$ and use the chart $x=o+\lambda Y$.
Everything remains strictly inside $(1/4,3/4)^3$.

Periodize the resulting compactly supported curl fields, loading during
the first unit interval and repeating thereafter. Compute the residual
at physical viscosity $\nu$, after the chart change. The fields and force
have zero mean and all the bounds in Lemma~\ref{bc:analytic}.
For a nonhalting run, loading has nonpositive logical height and each
later segment remains negative by the proposition. A terminal code has
physical height $1/2+\lambda$, which lies in the event. There is no
chart-boundary crossing that could create a spurious observation.
An initially terminal input is detected at the loading endpoint.
The initialized correspondence is exact; at checkpoints the number of
history records gives $n$, and the frontier's relative position together
with its absolute index $2+n$ also recovers the absolute work-head position.
\end{proof}

\subsection{Fixed columns and separate layers}\label{bc:columns-section}
\begin{proposition}[Column and layer transport]\label{bc:columns}
Let $S_i,T_i$ be finite rational planar rectangle families of positive side
lengths, each separately disjoint, with endpoints on the $B^{-2}$ grid,
where $B\ge2$ is an integer.
Let their assigned center-to-center positive reciprocal diagonal affine
maps satisfy $F_i(S_i)=T_i$. A finite effective
solenoidal field in $\R^3$ realizes all these maps on the rectangles in
$z=0$. Along each tracked path the first coordinate stays between its
initial and final values. The field has compact support and flat time
collars.
\end{proposition}
\begin{proof}\label{bc:columns-proof}
Give rectangle $i$ height $h_i=4i$. A planar cutoff equal to one near
its source and supported within padding $1/(4B^2)$ has support disjoint
from the other source cutoffs. Multiply it by a third-coordinate cutoff
which is one near the entire interval $[0,h_i]$. The curl of the potential
$(0,h_i x,0)$ times these cutoffs equals $(0,0,h_i)$ along that source
column. A common flat schedule lifts all rectangles. The corresponding
target cutoffs give a later descent by $-h_i$.

In the middle phase, let $p_i,q_i$ be the centers and $a_i>0$ the
first-coordinate scale. With phase parameter $\mu\in[0,1]$, set
$c_i=p_i+\mu(q_i-p_i)$, $l_i=1+\mu(a_i-1)$ and
\[
 H_i(\mu,x,y)=d_{i,1}y-d_{i,2}x+
       \frac{a_i-1}{l_i}(x-c_{i,1})(y-c_{i,2}),
 \qquad d_i=q_i-p_i.
\]
On a neighborhood of the desired path, the planar Hamiltonian field
$(\partial_yH_i,-\partial_xH_i)$ generates
$c_i+\diag(l_i,l_i^{-1})(\zeta-p_i)$.
The first coordinate is the convex interpolation of its endpoint values.
The second half-width obeys
\[
 \frac{r_{i,2}}{(1-\mu)+\mu a_i}
 \le (1-\mu)r_{i,2}+\mu r_{i,2}/a_i,
\]
by convexity of the reciprocal function. Consequently the swept rectangle
lies in the coordinate bounding rectangle of the two endpoints. Localize
$H_i$ around that rectangle in the layer $z=h_i$, with third padding less
than one, and take the curl of $(0,0,H_i)$ times that cutoff. For a
flat middle ramp $\mu=\eta_2(t)$, multiply the resulting phase field by
$\eta_2'(t)$. Layers are disjoint, so the sum has the required action. Lift and descent preserve
planar coordinates. The exact paths and ODE uniqueness establish the
full-rectangle and intermediate-coordinate assertions.
\end{proof}

\subsection{Target containers and prescribed affine images}
\label{bc:contained-targets}
The named targets in Propositions~\ref{bc:ordered}, \ref{bc:simultaneous}
and~\ref{bc:columns} can also be rational \emph{containers} for the actual
images. More precisely, replace the onto condition by
\[
 F_i(D_i)\subseteq S_i\quad\text{in Proposition~\ref{bc:ordered}},
 \qquad
 F_i(S_i)\subseteq T_i\quad\text{in the other two propositions}.
\]
Keep the other geometric hypotheses. For ordered transport retain the
center-to-center condition. For simultaneous and column transport the
actual endpoint center is $F_i(p_i)$, where $p_i$ is the source center;
it need not be the center of the container. The maps have the same positive
diagonal determinant-one linear parts as before, including the specified
three scale choices in Proposition~\ref{bc:simultaneous}.
All transport and path conclusions remain valid. For arbitrary real map
coefficients this is a smooth existence statement. The construction is
effective when those coefficients are computable, and otherwise its
evaluation is relative to them. An empty list again uses the zero field.

Here are the necessary modifications of the proofs. In ordered transport,
write $r_{i,j}$ for a source half-width, $w_{i,j}$ for the corresponding
container half-width, and $a_{i,j}>0$ for the diagonal factor. Containment
gives $a_{i,j}r_{i,j}\le w_{i,j}$, hence
\[
 r_{i,j}a_{i,j}^{\theta}
 \le\max(r_{i,j},a_{i,j}r_{i,j})
 \le\max(r_{i,j},w_{i,j})<L,\qquad 0\le\theta\le1.
\]
The rational source and container widths therefore bound every parked
deformation. After deformation, use translates of the rational target
containers as envelopes for the actual images. Deliver in increasing
order of the containers' right endpoints. If an earlier delivered image
meets the moving image in the last two projections, their containers meet
in those projections. Disjointness and the endpoint order force the earlier
container wholly to the left of the active target container. Its image
therefore misses the entire leftward sweep, which stops in that active
container. The original extraction and storage arguments are unchanged.
The same rational envelopes give positive cutoff margins. They do not
require the actual image endpoints to be rational. The central-cross-section
height guard is unchanged because its third offset remains zero.

For simultaneous transport, interpolate toward $F_i(p_i)$. The transported
source padding at descent has width at most $B\varepsilon$ outside its
actual image, so it lies inside the $B\varepsilon$ enlargement of the
target container. The original grid-gap and private-height estimates
still separate all supports. The first-coordinate convex identity uses
the actual value $F_i(y)$ and remains exact.

For column transport use the same actual endpoint center. The reciprocal
convexity estimate bounds each moving second half-width by the linear
interpolation of the source and actual-image half-widths. Thus the entire
middle sweep lies in the coordinate bounding rectangle of the source and
its rational target container. The target-column cutoff equals one on
the actual image, so the descent is unchanged. This also proves the stated
first-coordinate bound for each point. Finally, a positive computable
scale admits computable positive lower and finite upper bounds; its
logarithm, powers and the displayed reciprocal interpolations are
computable. The containment and determinant identities are hypotheses,
not tests performed by the construction.

\begin{theorem}[A half-space test by fixed columns]\label{bc:column-event}
On $\R^3$, the frontier recorder with odd digits in base $2m+1$ realizes
halting from the fixed label $0$ by the event $x_1<-1$. It gives a compactly
supported general force, periodic after time one, with zero initial
velocity and pressure and all mixed derivatives bounded.
\end{theorem}
\begin{proof}\label{bc:column-event-proof}
Put nonterminal state offsets at $3,6,\ldots$ and terminal offsets at
$-3,-6,\ldots$ in the first coordinate. Lemma~\ref{lem:bb-onecell} and
Proposition~\ref{bc:columns} give a whole-rectangle phase map. A localized
planar Hamiltonian translation loads the origin to the rational initial
code in $[0,1]$ time. A nonterminal loading segment has nonnegative first
coordinate, and every later nonterminal segment has first coordinate at
least three. Terminal codes have first coordinate below $-1$. This proves
both directions, including initial halting. Repeat the phase map and use
Lemma~\ref{bc:analytic}. Although this route only requires planar coding
rectangles, its cutoff plateaus also give a positive neighborhood on which
the same affine formulas hold.
\end{proof}

\subsection{Delaying the original head movement}\label{bc:delayed-column-section}
Use Lemma~\ref{lem:bb-direction} with singleton passive track removed.
Rename $(M(q,s),G(r),A(q,d),B(q,d))$ as
$(R_{q,s},F_r,C_{q,d},G_{q,d})$ and $(\bot,F)$ as $(\perp,*)$.
These bijections identify every row and nonfrontier guard on the full
domain, including all $q,d$ controls. The frontier starts at one, and the
work movement occurs only on return. At checkpoint $n$, write $i$ for
the work-head position and $j=n+1$ for the frontier position. The exact count is
$2(j-i)+3$: the right and left continuing loops have respectively
$j-i-1$ and $j-i$ instructions, with four other instructions.

\begin{theorem}[The delayed-head column realization]\label{bc:delayed-column-event}
The table of Lemma~\ref{lem:bb-direction} has a realization on $\R^3$
with fixed label $0$ and detector $x_1<-1$, compact spatial support,
period one after time one, bounded mixed derivatives and uniformly bounded
kinetic energy. The prescribed pressure is zero. Uniqueness holds in the
whole-space comparison class, in particular when the competitor and its
first spatial derivatives are bounded on every finite time interval.
\end{theorem}
\begin{proof}\label{bc:delayed-column-event-proof}
Assign odd digits $1,3,\ldots,2m-1$ in base $B=2m+2$, with state
offsets positive multiples of three except for negative terminal offsets.
The tail of the blank symbol contributes
$d_{\rm blank}B^{-k}/(B-1)$ after a finite prefix, so initialization is
rational even though the blank digit is nonzero. Full affine branches
follow from the lemma and Lemma~\ref{lem:bb-onecell}. Use the fixed-column
proof with heights $4i$ and the smaller padding $1/(10B^2)$; schedule a
slot $[a,b]$ by $\sigma(2(t-a)/(b-a)-1/2)$, which is flat on collars
of both endpoints. These choices preserve the same separation proof.
A planar Hamiltonian loading segment starts at zero. The convex first-
coordinate guard excludes the detector throughout every nonhalting
transition and throughout a nonhalting load; a terminal code lies below
$-1$. Finite checkpoint simulation gives the equivalence. The field has
one compact support and uniform pointwise bounds, hence uniformly bounded
kinetic energy. Lemma~\ref{bc:analytic} supplies all remaining assertions.
\end{proof}

\subsection{Recorded pointers and seven-slot storage}\label{bc:recorded-pointer-section}
\begin{theorem}[A fixed box test with recorded pointers]
\label{bc:recorded-pointer-event}
On $\R^3$, the fixed label $(2,0,0)$ and open box
$(0,1)^2\times(-1/2,1/2)$ realize halting by an effective general force
with compact spatial support, zero initial velocity and pressure, bounded
mixed derivatives, and period one after time one. The construction moves
whole solid boxes. Uniqueness holds among smooth competitors with
\[
 u\in C^1([0,T];L^2),\quad \nabla u,\Delta u\in C([0,T];L^2),
 \quad p,\nabla p\in C([0,T];L^2),\quad
 \sup_{[0,T]\times\R^3}(|u|+|\nabla u|)<\infty
\]
for every finite $T$.
\end{theorem}
\begin{proof}\label{bc:recorded-pointer-event-proof}
Normalize to one halt $h$. For $c=(q,b)\in\mathcal D=(Q\setminus\{h\})\times\Gamma$
write the instruction as $(q^+(c),b^+(c),m(c))$.
Use histories
\[
 \Lambda_0=\{\varnothing\}\sqcup\{H_e:e\in\mathcal D\sqcup\{\star\}\},
 \qquad \Lambda=\Lambda_0\sqcup\{P_e:e\in\mathcal D\sqcup\{\star\}\}.
\]
The full alphabet is $\mathcal A=\Gamma\times\Lambda\times\{0,1\}$,
whose tracks are data, history and return mark. Controls are
$S_q,L_q,M_c,R_c,F_c$, and the complete table is
\[
\begin{array}{c|c|c|l}
\text{input}&\text{output}&\text{move}&\text{restriction}\\\hline
S_q(b,l,0)&M_c(b^+(c),l,0)&m(c)&c=(q,b),\ l\in\Lambda_0\\
M_c(b,l,0)&R_c(b,l,1)&1&l\in\Lambda_0\\
R_c(b,l,0)&R_c(b,l,0)&1&l\in\Lambda_0\\
R_c(b,P_e,0)&F_c(b,H_e,0)&1&e\in\mathcal D\sqcup\{\star\}\\
F_c(b,\varnothing,0)&L_{q^+(c)}(b,P_c,0)&-1&\\
L_q(b,l,0)&L_q(b,l,0)&-1&l\in\Lambda_0\\
L_q(b,l,1)&S_q(b,l,0)&0&l\in\Lambda_0.
\end{array}
\]
Initially put $P_\star$ at absolute cell three, empty history elsewhere,
marks zero, and control $S_{q_0}$. At ready checkpoint $k$, the pointer
is at $r=3+k$, all cells to its right are empty, and the work configuration
is the original one. After the first row the work head satisfies
$j'\le k+1\le r-2$. The scan marks it, reaches the pointer, changes
$P_e$ into $H_e$, writes $P_c$ in the next empty cell, and returns to
erase the mark. Both scans are finite and positive. This proves the
initialized simulation. The initial pointer tag is retained as $H_\star$
after the first cycle, so the absolute cell-three reference is permanent.

The inverse works on all allowed tapes. Into $M_c$ the tag fixes the old
instruction and move; into $R_c$ mark one distinguishes entry from a loop;
into $F_c$ the written $H_e$ identifies the old pointer $P_e$; into $L_q$
a pointer output identifies entry, including the old $F_c$, whereas a loop
writes a nonpointer; into $S_q$ only mark erasure applies. Incoming shifts
are fixed in every target control. The remaining written components retain
all unchanged symbols, proving the two properties of
Lemma~\ref{lem:bb-onecell}.

Let $m_Q$ be the number of controls, and number them bijectively by
$\iota(s)\in\{0,\ldots,m_Q-1\}$ with $\iota(S_h)=0$.
Use odd digits in base $B=2|\mathcal A|+1$ and state offsets
$o_s=3\iota(s)$. The physical
code is $(o_s+L,R,0)$. Terminal codes lie strictly in $(0,1)^2\times\{0\}$;
nonterminal first coordinates are at least three. Form the full branch
boxes by taking the product with $[-1,1]$. Their planar sides have length
at most one. Choose $H_x=3m_Q+1$ and
parking centers $\pi_i=(H_x+10+4i,0)$, all at height zero.

Here is a storage construction distinct from ordered low extraction.
Evacuate each source in three successive operations: lift its center to
height ten, translate horizontally there to $\pi_i$, and lower it to
zero. Once every box is parked, process each in four operations: deform
its planar shape at the parking center, lift it to height ten, translate
it there to its target center, and lower it into the target. This uses
$7N$ slots for $N$ branches. During deformation use the positive scale
$l=1+\theta(\lambda_i-1)$ and its reciprocal, not an exponential scale.
Both side lengths remain bounded by the maximum of their endpoint lengths,
which is at most one. The third half-width remains one.

Let $\delta$ be the minimum of one and all positive source-source and
target-target planar distances, omitting empty sets of pairs. Use cutoff
plateaus with padding $\varepsilon=\delta/8$ and support padding
$2\varepsilon$. During a lift or descent the planar projection separates
the moving box from the current sources or delivered targets. Parked
boxes lie outside their entire first-coordinate range, and parking centers
have spacing four. During a horizontal translation the box lies in
$9\le z\le11$, whereas all stationary boxes lie in $-1\le z\le1$.
These positive gaps also separate the padded supports. Lemma~\ref{lem:bb-curl}
therefore realizes every operation without disturbing a stationary box.
The full assigned affine map follows by composition.

For a branch whose target control is nonterminal, every full point has
first coordinate at least three, including the storage and transfer paths.
Load the fixed label
$(2,0,0)$ to the rational initial code by a localized straight translation.
If the code is nonterminal, the loading segment has first coordinate at
least two; if it is terminal, the endpoint is in the observation box.
At later integer samples the exact recorder code gives terminal entry,
and the preceding bounds exclude the observation at all other times of a
nonhalting run. Repeat the finite transport field and apply
Lemma~\ref{bc:analytic}. Finally the displayed comparison class implies
$u\in C H^2$, since $\|u\|_{H^2}\le C(\|u\|_2+\|\Delta u\|_2)$
by the Fourier multiplier identity; its pressure assumptions give $p\in C H^1$.
Thus the stated uniqueness follows from the common whole-space comparison,
without changing this route's original hypotheses.
\end{proof}

\subsection{Planar detours behind a vertical observation guard}
\label{bc:guarded-planar-section}


\begin{theorem}[A guarded planar observation with unrestricted pressure]
\label{bc:guarded-planar-event}
On $\R^3$, the fixed label $(-4,0,0)$ and fixed box
$(-1,2)^2\times(-1,1)$ realize halting by a compactly supported general
force with zero initial velocity and prescribed pressure and all mixed
derivatives bounded. One choice is periodic after time one. A separate
choice satisfies $\|f(t)\|_\infty=O((1+t)^{-2/3})$ and
$f\in L^2([0,\infty)\times\R^3)$.
Uniqueness holds among smooth competitors with
$u\in C([0,T];H^2)\cap C^1([0,T];L^2)$ and bounded $u,\nabla u$
on each finite interval, without any spatial integrability condition on
pressure. Pressure is unique modulo a function of time.
\end{theorem}
\begin{proof}\label{bc:guarded-planar-event-proof}
Use the frontier recorder after normalization to one halting state, and
write $\mathcal A$ for its full cell alphabet and $m_Q$ for its number
of controls. Number the controls bijectively by
$\iota(s)\in\{0,\ldots,m_Q-1\}$, giving the terminal control index zero.
Use odd base $B=2|\mathcal A|+1$ and first-coordinate offsets
$o_s=4\iota(s)$. Code the
configuration by $(o_s+L,R)$. The origin of the original tape can be
recovered from the number of history records and the frontier position,
so this head-relative coding loses no initialized information.
The full planar branches have reciprocal positive diagonal maps.

We first realize finite separately disjoint source and target rectangle
families $S_i,T_i$ with rational endpoints and positive side lengths,
equipped with assigned positive
reciprocal diagonal affine maps $F_i$ satisfying $F_i(S_i)=T_i$, by
compactly supported planar Hamiltonian motion. Write $c_i,q_i$ for the
source and target centers. With no
branches use zero. With one branch, interpolate its center directly and
use a positive linear first-coordinate scale and its reciprocal; there
are no stationary obstacles. For $N\ge2$, let $R=1+$ the largest endpoint
half-width and let $m_0>0$ be the minimum sup-norm separation within the
two center families. Put $K=1+32R/m_0$. Expand all source centers by
factors from one to $K$, leaving their shapes unchanged. Let
$M=\max_i\{|(Kc_i)_1|,|(Kq_i)_1|\}$ and choose parking centers
$P_i=(M+32Ri,0)$. Each union of expanded endpoint centers and parking
centers has separation at least $32R$.

Move expanded sources successively to their parking centers, deform them
there using the linear reciprocal scale, move them successively to the
expanded target centers, then contract those centers to their targets.
For each successive transfer put a square of sup-radius $4R$ around
every stationary center. Replace each nontrivial segment intersection
with an obstacle square by a counterclockwise path along its boundary.
The intersection points and corners are rational; the squares are disjoint,
so the resulting finite polygon stays at distance at least $4R$ from
all stationary centers. Flat schedules on its segments give smooth center
motion. The simultaneous expansion and contraction preserve a positive
coordinate gap for every pair, exactly as in the solid-box argument.
If $g_0,g_1$ are their minimum positive gaps, choose padding
$\varepsilon=\min(g_0,g_1,2R)/4$ for these phases. During individual
transfers use padding below $R$; the stationary half-widths are below $R$.
Thus all required cutoff supports avoid stationary rectangles.

For a moving center $c$ and horizontal scale $l$, the potential
\[
 H=\chi\left[\dot c_x(y-c_y)-\dot c_y(x-c_x)
                    +\frac{\dot l}{l}(x-c_x)(y-c_y)\right]
\]
gives the required motion on the cutoff plateau. The other current
rectangles are fixed. This proves every full planar branch, including
its boundary, by induction through the operations. A one-square loading
translation with half-width $1/4$ takes $(-4,0)$ to the initial rational
code. All these supports lie in a square $[-L,L]^2$, where one can choose
$L=10+S_0+4m_Q$ with $S_0$ a finite bound for routing vertices and cutoff
widths. The same $L$ covers every input of a fixed machine, since initial
codes are uniformly bounded.

The planar detours need not avoid the horizontal observation box.
We instead make them at height three. Choose a switch $\eta$ with
collars and set
$h(s)=3\{\eta(3s)-\eta(3s-2)\}$ on a unit interval. It lifts in the
first third, stays at three in the middle, and lowers in the last third.
Choose $\Gamma$ equal to one near $[-L,L]^2\times[-1,4]$, with compact
support, and $\beta(z)$ compactly supported and equal to one near three.
For any one of the planar fields $w$, zero outside its phase interval,
put
\[
 Z(s,x)=h'(s)\nabla\times(0,x_1\Gamma(x),0),\qquad
 V_w(s,x)=Z(s,x)+3\beta(x_3)(w(3s-1,x_1,x_2),0).
\]
Along each tracked path the first term is exact vertical translation.
The second term acts only in the middle third, at height three, and runs
the full planar phase. It is divergence free because $w$ is planar
solenoidal and $\beta$ is independent of the planar coordinates.
Thus the trajectory lifts at its exact source, performs every detour
outside the vertical detector interval, and lowers at its exact target.
Use this construction for the loader and for each repeated phase. While
its height can lie in $(-1,1)$, a nonterminal computational particle has
first coordinate at least four; the initial loading source has first
coordinate $-4$. A terminal code lies in $(0,1)^2\times\{0\}$.
These facts prove the all-real-time observation equivalence, including
an initially terminal input.

For the comparison assertion, the constructed field has compact support
and bounded mixed derivatives, so it satisfies the reference hypotheses
of Lemma~\ref{lem:b-comparison}. The displayed competitor class is exactly
case (v) of that lemma. Its proof obtains $\nabla p\in L^2$ from the
equation and cancels the pressure term by Lemma~\ref{lem:b-gradient}; it
does not require an integrable pressure representative. Hence the velocity
is unique, and pressure is determined up to a function of time.

Finally let $U(s)$ denote the eventually periodic field just constructed
and set $\rho(t)=(1+t)^{1/3}-1$, $\widehat U=\rho' U(\rho(t))$.
The new residual is
\[
 \widehat f=\rho''U+(\rho')^2\{U_s+(U\cdot\nabla)U\}
                          -\nu\rho'\Delta U,
 \quad \rho'=\tfrac13(1+t)^{-2/3},\quad
 \rho''=-\tfrac29(1+t)^{-5/3}.
\]
Uniform derivative bounds and the fixed compact support give the stated
decay and space-time $L^2$ estimate. Chain differentiation also gives all
mixed bounds. The clock is effective on $[0,\infty)$, starts at zero,
is finite at finite time and is onto; hence it preserves precisely the
same material event. The initial zero collar is preserved as well.
This slower force is an alternative to, not a periodic version of, the
first force.
\end{proof}
